
\textbf{Metodo gráfico.} Consiste en graficar la actividad contra el tiempo tomando en cuenta que despues de cada vida media la actividad original va decayendo a la mitad. (Se utiliza papel similogaritmico para graficarlo), siendo la grafica exclusiva para el Ir-192.

\begin{center}
  \includegraphics[scale=0.5]{graficaIridio.png}
\end{center}

% Empieza la sección de "Interacción de la Radiación con la Materia"
\chapter{Interacción de la Radiación con la Materia}

Se entiende por radiación ionizante, aquellas ondas electromagnéticas o partículas capaces de interactuar con los electrones de los átomos que componen la materia, para producir iones.

% Explicación sobre los iones
Los iones son átomos o moléculas cargados eléctricamente.

% Descripción de las ondas electromagnéticas que pueden producir ionización
Las ondas electromagnéticas capaces de producir ionización, son: los rayos X y gamma. Las partículas que producen ionización son: las alfa, beta, protones y neutrones.

% Definición de ionización
La ionización es el proceso por medio del cual un átomo o molécula adquiere una carga eléctrica, positiva o negativa, mediante la extracción o adición de electrones.

% Definición de excitación
Excitación es el estado donde se encuentra un excedente de energía que se elimina mediante la emisión de rayos gamma.

